\chapter{Sty (Hordeolum)}

\section*{Definition}
A sty (hordeolum) is an acute, painful infection of the eyelid margin involving the sebaceous glands of Zeis (external hordeolum) or meibomian glands (internal hordeolum). It is most commonly caused by Staphylococcus aureus.

\section*{What Happens in Your Body}
Bacterial infection (usually S. aureus) of the eyelid glands causes localized inflammation, pus formation, and swelling. An external hordeolum involves the glands of Zeis at the base of the eyelash follicle. An internal hordeolum involves the meibomian glands within the eyelid tarsal plate and is usually more painful. Inflammation leads to a tender, red, swollen nodule.

\section*{Causes}
\begin{itemize}
\item Bacterial infection: Staphylococcus aureus (most common)
\item Blepharitis (chronic eyelid inflammation is a risk factor)
\item Poor eyelid hygiene
\item Contact lens wear (improper hygiene)
\item Rosacea
\item Diabetes mellitus
\item Immunosuppression
\item Stress and fatigue
\end{itemize}

\section*{When to See a Doctor}
\begin{itemize}
\item Red, swollen, painful lump on the eyelid margin
\item Localized tenderness
\item Swelling of the entire eyelid (if severe)
\item Gritty or foreign body sensation
\item Watering of the eye
\item Crusting along the eyelashes
\item Pus may drain from the sty
\end{itemize}

\section*{Related Symptoms}
\begin{itemize}
\item Blepharitis (eyelid inflammation)
\item Chalazion (chronic, non-infectious meibomian cyst)
\item Conjunctivitis (pink eye)
\item Periorbital cellulitis (severe infection)
\end{itemize}

\section*{Self-Care / Home Management}
\begin{itemize}
\item Apply warm compresses for 10-15 minutes, 3-4 times daily
\item Gently massage the area after warm compresses
\item Do NOT squeeze or pop the sty
\item Clean the eyelid with diluted baby shampoo
\item Avoid eye makeup and contact lens wear until healed
\item Wash hands before and after touching the area
\end{itemize}

\section*{How Your Doctor Will Evaluate You}
\begin{itemize}
\item Clinical examination is diagnostic
\item Slit-lamp examination of the eyelid
\end{itemize}

\section*{Treatment Options}
\begin{itemize}
\item Warm compresses are the mainstay of treatment
\item Topical antibiotic ointments (erythromycin, bacitracin)
\item Topical antibiotic-steroid combination for inflammation
\item Oral antibiotics for recurrent or severe cases
\item Incision and drainage for large or persistent sties
\item Treat underlying blepharitis or rosacea
\end{itemize}

\section*{References}
\begin{thebibliography}{99}
\bibitem{sty:ref1} Lindsley K, Nichols JJ, Dickersin K. "Interventions for acute internal hordeolum." \textit{Cochrane Database of Systematic Reviews}, 2010;2:CD007742.
\bibitem{sty:ref2} Skorin L Jr. "Hordeolum and chalazion: Diagnosis and management." \textit{Optometry Clinics}, 1995;4(3):79--84.
\bibitem{sty:ref3} Gilchrist H, Lee GA, Lee G, et al. "Management of chalazia and hordeola." \textit{Australian Family Physician}, 2006;35(5):323--326.
\bibitem{sty:ref4} Patel MR, Quickert MH. "Hordeolum and chalazion: Treatment and prevention." \textit{Optometry}, 2002;73(10):599--604.
\bibitem{sty:ref5} Kim JH, Kim JS, Kim WH, et al. "Clinical features and treatment outcomes of hordeolum in adults." \textit{Korean Journal of Ophthalmology}, 2019;33(1):39--44.
\bibitem{sty:ref6} Tsai CC, Kau HC, Kao SC, et al. "Efficacy of probing the nasolacrimal duct with adjunctive intubation for congenital nasolacrimal duct obstruction." \textit{Journal of Pediatric Ophthalmology \& Strabismus}, 2004;41(5):274--279.
\end{thebibliography}
